% !PRINTED 592
\weberpage{592}
If one seeks the greatest common divisor of $\boldsymbol{A}_\mu$ and
$\boldsymbol{A}_{\mu_1}$, one obtains an integral function $\boldsymbol{A}_{\mathfrak{m}}$ whose coefficients % [source sic: the printed German singular Koeffizient conflicts with the plural predicate]
are integers in $\mathfrak{L}$ and whose roots are the quantities (9); and
if one finally frees $\boldsymbol{A}_{\mathfrak{m}}$ of the factors that it has in common
with any $\boldsymbol{A}_{\mathfrak{m}'}$ in which $\mathfrak{m}'$ is a proper divisor of $\mathfrak{m}$, one
obtains an integral function:
\weq{10}{T_{\mathfrak{m}}(x) = x^{\nu} - \tau_{\nu-1}x^{\nu-1} + \tau_{\nu-2}x^{\nu-2} - \cdots \pm \tau_0,}
whose roots are only those among the quantities (8) in which $\xi$
is relatively prime to $\mathfrak{m}$.
% [source sic: the roots have just been identified as the quantities (9), but this sentence calls them quantities (8)]

The degree $\nu$ of this function is to be determined according to Bd.\ II, \S\,168.
We shall call it the \gesperrt{ideal division function}. Its coefficients
are \gesperrt{integers} of the field $\mathfrak{L}$, and in
particular
\weq{11}{\tau_0 = \overset{\xi}{\Pi}\,\boldsymbol{s}\!\left(\frac{2\alpha\xi}{\mu}\right),}
in which $\xi$ runs through a complete system of numbers incongruent to one
another and coprime to $\mathfrak{m}$.

The roots of $T_{\mathfrak{m}}(x)$ are equal and opposite, and
we obtain an equation of degree $1/2\,\nu$ for the quantity
\[
\boldsymbol{s}\!\left(\frac{2\alpha\xi}{\mu}\right)^{2}.
\]
By adjunction of these quantities to the field $\mathfrak{L}$ there arises
a field $\mathfrak{T}_{\mathfrak{m}}$, which we likewise call a division field. It
results from the division field of \S\,154 by adjunction
of $x$. The group of this field is contained in the group of the numbers
$\xi\ (\mathrm{mod}\ \mathfrak{m})$, as follows easily from the multiplication formulae.

We note further that the general division field
$\mathfrak{T}_{\mathfrak{m}}$ can be reduced, by application of the addition theorem, to
the case where $\mathfrak{m}$ is a power of a prime ideal in $\Omega$.

\begin{center}
\textbf{\S\ 159. Prime Ideals.}
\end{center}

6. \gesperrt{If $\eta$ is a fractional number of the field $\Omega$ which in
reduced form has the odd ideal denominator $\mathfrak{a}$, and if $\mathfrak{m}$
is an ideal relatively prime to $\mathfrak{a}$, then the integer $\boldsymbol{s}(\eta)$ is relatively
prime to $\mathfrak{m}$.}

For let there be taken in $\Omega$ an integer $\mu$ divisible by $\mathfrak{a}$ and
 coprime to $2\mathfrak{m}$. Then $\mu\eta$ is an integer. In the
product \S\,157, (23)
%CONT>>
